% Einheit 3 von 4 – Klammern auflösen (bank/terme/e3.jsonl, zusammenbau v0.1)
%% TODO Zweigzeile Teil 1: was hier gelernt wird (Satz in Schülersprache aus dem Einheitstitel) – Einheit 3
\einheitenkopf[e3]{Einheit 3 von 4 · Klammern auflösen}
\zweigzeile{ab Kl. 7, je nach Buch bis Kl. 8 · keine P10-Aufgabe}
\setcounter{aufgabe}{16}

%% TODO Titel: Ich-kann-Satz fehlt in der Bank (Platzhalter: Kettenname) – Nr. 17
%% TODO Anweisung: ein Satz über der ersten Teilaufgabe fehlt in der Bank – Nr. 17
\begin{aufgabe}{Zeichen vor der Klammer feststellen}
\begin{teile}
\teil $7 + (a - 2)$ – was steht vor der Klammer? \leerfeld
\end{teile}
\end{aufgabe}

%% TODO Titel: Ich-kann-Satz fehlt in der Bank (Platzhalter: Kettenname) – Nr. 18
%% TODO Anweisung: ein Satz über der ersten Teilaufgabe fehlt in der Bank – Nr. 18
\begin{aufgabe}{Klammern}
%% TODO \rechenplatz{n}: Schrittzahl des Grundfalls fehlt in der Bank (nur bei fester Schreibform, 2.3 b)
\begin{gleichungsraster}[2]
\gl{\text{Löse die Klammer auf: $4x + (3y - 7)$}} &
\end{gleichungsraster}
\begin{teile}
\teil Löse die Klammer mit der Tabelle auf: $5 \cdot (x + 3)$

\sachtabelle{ccc}{$\cdot$ & $x$ & $+3$}{$5$ & \leerzelle & \leerzelle}
\leerfeld
\end{teile}
\begin{gleichungsraster}[2]
\gl{\text{Löse die Klammer auf: $5a - (2b + 4)$}} &
\gl{\text{Löse die Klammer auf: $7a - (2b - 5 + c)$}} \\
\gl{\text{Löse die Klammer auf: $-2 \cdot (x + 5)$}} &
\gl{\text{Löse die Klammer auf: $5x + 2 \cdot (x + 6)$}} \\
\gl{\text{Löse die Klammer auf: $3 \cdot (2x + 5) - 2 \cdot (x - 4)$}} & \\
\end{gleichungsraster}
\end{aufgabe}

%% TODO Titel: Ich-kann-Satz fehlt in der Bank (Platzhalter: Kettenname) – Nr. 19
%% TODO Anweisung: ein Satz über der ersten Teilaufgabe fehlt in der Bank – Nr. 19
\begin{aufgabe}{Klammern – Fehler finden, Begründen, Darstellungswechsel}
\begin{teile}
\teil Finde den Fehler und rechne richtig: \rechnung{12 - (a - 5) &= 12 - a - 5 \\ &= 7 - a}
\teil Sind $5 - (x - 2)$ und $7 - x$ gleichwertig? Begründe.
\teil Flächeninhalt des Rechtecks (Maße in cm) – Term mit Klammer und ohne Klammer?

\viereck[seiten={x+2,5,x+2,5}]{(0,0)}{(6,0)}{(6,2.5)}{(0,2.5)}
\leerfeld
\end{teile}
\end{aufgabe}
